\section{Memory access characteristics}
\label{section:core:memory-access-characteristics}


\paragraph{Hypothesis}

If memory bandwidth used is low and the MFlops rate is low, too, we often have a
situation where the code issues many scattered memory accesses all over the
place all the time.
Every time we miss a cache, we have to access the main memory and pay for the
latency. 
We will never be able to saturate its theoretical bandwidth even though we
access the memory all the time.
Along the same lines, we will never be able to exploit our compute capabilities,
as we never get the data into the registers quickly enough.



\begin{assessmentcrime}
 80\% of the runtime might be spent on 20\% of the code.
 For the performance analysis, it is not appropriate to focus on this 20\% of
 the code only.
 Instead, you have to focus on all code that does not do any science if this
 code share is too big.
\end{assessmentcrime}

Too big means more than 20\%. Flame graphs helfen also gar net. Sie sind sogar
counter-productive bzgl Outcome. 
Wie passt die Aussage eben zu der Tatsache, dass wir nix ueber den Code wissen
sollten.
Developer muss Black/white List schicken.  Aber eben net ``was wird gebraucht''
sondern nur ``wo wird gerechnet''.



